\new{The core concept of \rap is to incorporate a penalty for the model’s performance based on the degree of repetition in the generated response. To achieve this, we propose an efficient algorithm, \alg, for detecting and quantifying repetition in text. We introduce the notion of the \textit{repetition ratio} (\rr), which measures the proportion of repeated content. Additionally, we define \textit{repetition-aware performance} (\rap), a metric that integrates repetition penalties into standard performance evaluations, enabling tuning of the repetition penalty parameter.}
% \new{The main idea of \ours is to penalize model's performance based on the level of repetition in the response. To this end, we propose an efficient algorithm to detect and measure repetition in a generated text. We introduce the notion of \textit{repetition ratio} (\rr) which captures the proportion of repeated content, and define \textit{repetition-aware performance} (RAP) which quantitatively integrates repetition penalties into performance evaluation metrics. }
% In this section, we define the problem of repetition detection. \new{We introduce the notion of \textit{repetition ratio} (\rr) which captures the proportion of repeated content.} 
% We present our method for detecting repetition \new{and define \textit{repetition-aware performance} (RAP)} which quantitatively integrates repetition penalties into performance evaluation metrics. 
% This facilitates identifying the trade-off between repetition and task performance.
%
\new{Given a text sequence generated by an LLM, we first detect \textit{non-word characters} (\nwc) repetition, that includes white spaces and non-alphanumeric characters such as punctuation and symbols.}
A sequence has \nwc repetition if it contains at least \textit{k} consecutive identical \nwc. After carefully checking the outputs of various LLMs, we choose $k=5$ to ignore cases where \nwc are used for formatting the output.
We then remove these repeated \nwc and detect \textit{text repetition}. 
% For text repetition detection, we assume all \nwc have been removed and use the remaining content for detection.

\subsection{Definitions}
% Definition of REPETITION. Types of repetition: new lines and text. Can borrow the definition of the other paper. Define ``new lines'' repetition (why we need this type of repetition)

% Define the main problem (find a name for this problem): to find the best repetition parameter (trade-off) that can obtain good performance and suppress the repetition issue. 

% Solution: Repetition factor is solution 

% Propose a framework to quantify the repetition behaviour which is currently not addressed. balance between performance and repetition. 

Repetition in LLM-generated content can occur not only in sequences of words separated by spaces but also within a single, super-long ``word'' that contains no spaces. Both of these issues are equally important. Therefore, we extend the definition of text repetition from \citet{Fu2020ATA} to cover not only the repetition of word sequences but also the repetition of characters within no-space content. 
% , i.e., content witout spaces.
%
We denote an LLM's generated sequence as $S=[w_1, w_2, \cdots, w_{|S|}]$ where $w_i$ is the $i^{th}$ character in $S$ and $|S|$ is the length of the sequence in terms of characters. We denote $S_{p:q}=[w_p, w_{p+1}, \cdots, w_{q}]$, where $1\leq p < q \leq |S|$, as a continuous subsequence of $S$. 
% We now define the notion of \textit{equal subsequences}. 

\begin{theorem}[Equal Subsequences]
\label{def:equal_sub}
Subsequences $S_{a:b}$ and $S_{c:d}$ are said to be equal if $b-a = d-c$ and $w_{a+i}=w_{c+i}$, for $\forall i \in [0, b-a]$. 
\end{theorem}

% Following \citet{Fu2020ATA}, we define that a subsequence is repeated if it equals to its consecutive subsequence.
\begin{theorem}[Repeated Subsequence]
\label{def:rep_sub}
Given a sequence $S=[w_1, w_2, \cdots, w_{|S|}]$, a subsequence $S_{p:q}$, $1 \leq p < q \leq \frac{|S|+p-1}{2}$ is a repeated subsequence if it equals to its consecutive subsequence, $S_{q+1:2q-p+1}$.
% A subsequence $s_{p:q}$ is a repeated subsequence if it equals to its consecutive subsequence $s_{q+1:2q-p+1}$.
\end{theorem}

\begin{theorem}[Maximal Repeated Subsequence]
\label{def:rep_sub}
A subsequence $S_{p:q}$ is said to be a maximal repeated subsequence if $\nexists c, d$ such that $S_{c:d}$ is a repeated subsequence, where $c \leq p < q \leq d$ and $d-c > q-p$.
% Given a sequence $S=[w_1, w_2, \cdots, w_{|S|}]$, a subsequence $s_{p:q}$, $1 \leq p < q \leq \frac{|S|+p-1}{2}$ is a repeated subsequence if $s_{p:q}$ and $s_{q+1:2q-p+1}$ are equal.
\end{theorem}

% A \textit{maximal repeated subsequence} (\mrs) in a sequence is the longest subsequence that repeats without interruption, capturing the largest repeated unit. The repetition length of a \mrs is measured by the length of the repeated content for that pattern, ignoring the first occurence of the pattern. 

% \noindent
A \textit{maximal repeated subsequence} (\mrs) in a sequence refers to the longest  subsequence that occurs repeatedly without interruption, representing the largest repeating unit. There may be multiple \mrses within a text, each representing a distinct repeated pattern.  \new{The repetition length of an \mrs is defined as the length of the repeated content, excluding the first occurrence of the pattern. The repetition length of the entire sequence is calculated as the sum of the repetition lengths of all the \mrses within the sequence}.

% A maximal repeated subsequence (\mrs) in a text sequence is the longest consecutive subsequence that repeats as a whole without any interruption. \mrs encompasses the entire repeated sequence, capturing the largest possible repeated unit. There may be multiple \mrs within a text sequence, each representing a distinct longest repeated pattern. For example, in the generated text shown in Figure \ref{fig:overall}, the MRS are ``apple apple orange'' and ``tree tree cup tree cup''. We now define the \textit{text repetition detection} problem.

\begin{theorem}[Text Repetition Detection Problem]
\label{def:rep_sub}
Given a sequence $S$, identify all maximal repeated subsequences (\mrses) within $S$ and determine their corresponding repetition lengths.
\end{theorem}

\new{\begin{theorem}[Repetition Ratio -- \rr]
\label{def:rep_sub}
The \textit{repetition ratio} (\rr) of a sequence $S$ is defined as the proportion of repeated content relative to the length of the sequence: 
% find all maximal repeated subsequences in $s$ and the number of times each of such the sequences repeated. 
\begin{align}
\label{eq:rr}
% \text{\rr}(S) = \frac{\sum_{S_i \in R_S}|S_i| \times (\text{freq(}S_i\text{)}-1)}{|S|}  
\text{\rr}(S) = \frac{R(S)}{|S|}  
\end{align}
where $R(S)$ is the repetition length of the entire sequence $S$. 
% the set of all \mrses in the input sequence $S$, \textit{freq($S_i$)} is the frequency of \mrs $S_i$. 
\end{theorem}
% The \rr of a set of responses $D$ can be computed as the mean of the \rr of the responses in $D$, but that will fail to differentiate cases where the response length   is computed as follows.
% \begin{align}
%     \label{eq:rr_set}
%     \text{\rr}(D) = _{S\in D}
% \end{align}
% }
}

% \noindent
In Figure \ref{fig:overall}, the repeated content includes \nwc repetition (repeated dots), and text repetition (``United States Dollar''). The \rr computed using Equation \ref{eq:rr} is 0.373, i.e., about 37\% of the text is repeated content. 

% We now define the text repetition detection problem.
% The text repetition detection task involves, first identifying each \mrs $s_{p:q}$, $1\leq p < q < |S|$, and secondly, counting the frequency of each \mrs. Back the example in Figure \ref{fig:overall}, the \mrs ``apple apple orange'' occurs 3 times and the \mrs ``tree tree cup tree cup'' occurs 2 times.


\noindent
\textbf{Repetition Penalty Parameter (\rp).} \rp is applied during token sampling to address repetition in generated text \cite{Keskar2019CTRLAC}. Specifically, the probability of predicting the $i^{th}$ token is computed as:
\begin{align}
    p_i=\frac{\text{exp}(x_i/(T\cdot I(i\in g)))}{\sum_j\text{exp}(x_j/T\cdot I(j\in g)))}
\end{align}
where $I(c) = \theta$ if $c$ is True else 1, $T$ is the temperature value, $g$ is the list of generated tokens, and $\theta$ is the \rp \cite{Keskar2019CTRLAC}. \rp adjusts the probabilities of previously generated tokens, penalizing them to reduce their likelihood of being selected again. The extent of this adjustment is governed by the \rp value: higher \rp values impose stronger penalties, effectively suppressing repetitive outputs. However, this comes with a trade-off, as overly penalizing tokens can potentially impact the coherence and fluency of the generated text. We now present the repetition detection algorithm and \rap (repetition-aware performance), which facilitate tuning \rp to achieve an optimal balance. 



\subsection{Repetition Detection Algorithm (\alg)}
\label{sec:rep_det}
% We meticulously designed and implemented an efficient algorithm, namely \alg, which employs regular expression to detect repetitions for both \whitespace repetition and text repetition. Our algorithm can also detect repetitions within a single no-space content, an issue that is not commonly covered \cite{Fu2020ATA} . 
% \alg is shown in Algorithm \ref{alg:detect_repetitions}. 
% \alg uses two regular expression patterns: one for detecting sequences of \whitespace and another for identifying repeated text sequences. 
% It first detects all \whitespace repetition, records the repetition length, and remove all \whitespace repetition in the input text for the next step of text repetition detection (Lines 6--7). 
% In text repetition detection, we adopt the second pattern to identify all \mrs and the corresponding repetition length (Line 9--14). 
% Ultimately, \alg returns the lengths of the \whitespace repetition, text repetition, and the total repetitions, which is the sum of both \whitespace and text repetitions.
\alg adopts \textit{regular expressions} to detect both \nwc and text repetitions. \alg can detect repetitions within single no-space content, an issue often overlooked \cite{Fu2020ATA}. 
% \alg employs two regular expression patterns for detecting \nwc repetiton and text repetition. 
The regular expression for detecting \nwc repetition is as follows:
\begin{align}
    [\textbackslash s\textbackslash W]\{5,\}
\end{align}
\noindent 
which matches any sequence of 5 or more consecutive characters where each character is either a whitespace or a \nwc. The regular expression for identifying text repetition is: 
\begin{align}
 (?P<r>.\{5\}.*?)(?:[\textbackslash s\textbackslash W]*(?P=r)+   
\end{align}
which matches a sequence of at least 5 alphanumeric (captured in the group $r$), followed by one or more repetitions of the same sequence, potentially separated by whitespace or non-word characters.
% : one for \nwc sequences and another for repeated text sequences. The pattern for detecting \nwc repetition is ``\text{([\\s\\W]{5,})}''. The pattern for detecting text repetition is ``(?P<r>.\{5\}.*?)\\\hspace{2.5cm}(?:[\textbackslash s\textbackslash W]*(?P=r)+''. 
\alg first detects and records all \nwc repetitions, removes them, and then identifies text repetitions. For text repetition, it identifies all \mrses and their frequencies. Finally, \alg returns the repetition lengths of \nwc and text repetition. \alg is detailed in Algorithm \ref{alg:detect_repetitions}.
% (Appendix).

\begin{algorithm}[ht]
\small
\caption{Repetition Detection (\alg)}
% updated by Donghao on Oct 9
\label{alg:detect_repetitions}
\begin{algorithmic}[1]
\State \textbf{Input:} $G_{text}$ \Comment{Input text}
\State \textbf{Output:} RL$_{\nwc}$, RL$_{text}$, RL$_{total}$ \Comment{Repetition length of \nwc, text, and total, respectively}
\State P$_{\nwc}$ = \texttt{re.compile(r"([\textbackslash s\textbackslash W]\{5,\})")}  \Comment{Regex pattern for detecting \nwc repetition}
\State P$_{text}$ = \texttt{re.compile(r"(?P<r>.\{5\}.*?)\\\hspace{2.5cm}(?:[\textbackslash s\textbackslash W]*(?P=r)+")} \Comment{Regex pattern for detecting text repetition}\\
\Comment{\textbf{STEP 1. Detecting \nwc repetition}}
    \State $G_{text}'$ = P$_{\nwc}$.sub("\textbackslash t
", $G_{text}$)  \Comment{Remove \nwc repetition}
    \State RL$_{\nwc}$ = len($G_{text}$) - len($G'_{text}$) \Comment{Compute \nwc repetition} \\
\Comment{\textbf{STEP 2. Detecting text repetition}}
    \State RL$_{text}$ = 0
    \State matches = P$_{text}$.finditer($G_{text}'$)
    \For{each match in matches}
        \State (start, end) = match.span()
        \State RL$_{text}$ += end - start - len(match.group(1))  \Comment{Aggregate text repetition lengths}
    \EndFor\\
    RL$_{total}$ = RL$_{\nwc}$ + RL$_{text}$
    \State \Return RL$_{\nwc}$, RL$_{text}$, RL$_{total}$

% \Function{detect\_repetitions}{text}
%     \State text, a $\gets$ del\_excessive\_nwcs(text)
%     \State b $\gets$ detect\_text\_repetitions(text)
%     \State 
    % \Return (a, b, a + b)
% \EndFunction
\end{algorithmic}
\end{algorithm}



% including the length and frequency of each repeated pattern. 


% \subsubsection{Baseline Method \todo{to shorten and move to experimental setting}}
% The baseline method based on \cite{Fu2020ATA} detects the repetition within a text by detecting same word phrases that are adjacent to each other. In order to simulate this, we take the tokenized words of the input text and store them in a suffix array \cite{10.5555/320176.320218} as suffixes of the text. Here, we define the suffix of a text as a word phrase (substring) that occurs at the end of a text. For example, a text "This is repeated this is repeated text" has 7 suffixes in the following order: 0. "This is repeated this is repeated text" (the original sentence), 1. "is repeated this is repeated text", 2. "repeated this is repeated text", 3. "this is repeated text", 4. "is repeated text", 5. "repeated text", and 6. "text". Suffix array stores these phrases in lexicographical order. Hence after sorting, the suffix array stores the phrase indexes of this example as 4, 1, 5, 2, 6, 3, 0. Figure \ref{fig:suffix_array_simulation} shows the structure of the suffix array example after sorting.  

% Once we have the suffix array, we can extract all possible repeated phrases by scanning through the array. There are $N$ passes. For each pass, it goes through each row in the array, takes $N$ token(s) as a word phrase, and compares it with the next $N$ token(s) in the suffix array. These two word phrases are considered repeating if (1). they are the same word phrase and (2). the distance of the suffix array indexes is equal to the length of the phrase. Figure \ref{fig:suffix_array_simulation} simulates the first three passes of "This is repeated this is repeated text" example.  

% After we extract all possible repeated phrases, we use greedy approach to extract the longest combination of repeated phrases from the provided list. 

% The final complexity of this algorithm is at $O(N^2)$ where $N$ is the size of the tokenized words. % this depends on the greedy approach

% WARNING: still a draft, need verification against the original code one more time if we're going to use this / put this in appendix.

% \begin{algorithm}
% \caption{Get adjacent repeated phrases using suffix array}\label{alg:cap}
% \begin{algorithmic}
% \State $w \gets \text{tokenize}(text)$
% \State $sa \gets \text{build\_suffix\_array}(w) $
% \State $sal \gets \text{[]}$ \Comment{\textcolor{blue}{suffix array's label}}
% \For{$i$ in 0 ... len($w$)}
%     \State $sal[i] \gets w[i:\text{len}(w)]$
% \EndFor
% \\
% \State $arp \gets \text{[]}$
% \For{$c$ in 0 ... len{}($w$)}:
%     \State $rp \gets \text{[]}$
%     \For{$r$ in 1 ... len{}($sa$)}:
%         \If{$sal[sa[r]][c]$ = $sal[sa[r-1]][c]$}
%             \State $rp[\text{len}(rp)-1].\text{append}(sa[r]) $
%         \Else{}
%             \State $rp.\text{append}([sa[r]])$
%         \EndIf
%     \EndFor

%     \State $rp \gets \text{get\_adjacent\_phrase}(rp, c+1)$           
%     \State $arp\text{.extend}(rp)$
% \EndFor
% \\
% \State \text{return} $arp$
% \end{algorithmic}
% \end{algorithm}


\subsection{Repetition-Aware Performance (\rap)}
\label{sec:rap}
% Using \alg, we measure the repeated length of each LLM response. 
\new{After measuring the repeated length of each LLM response, we compute the \rr for the entire dataset based on Equation~\ref{eq:rr}:
%to assess the repetition-aware performance (\rap):
 %A simple approach is to average the \rr of each response. However, this method fails to capture extreme cases where responses are excessively long and highly repetitive. Inspired by the F1 score, which is computed using averaged precision and recall rather than averaging F1 scores, we propose the following formula to compute \rr for the entire dataset:
\begin{align}
\label{eq:rr_dataset}
    \text{\rr}_D = \frac{\sum_{S\in D} R(S)}{\sum_{S\in D} |S|}
\end{align}
where $\text{\rr}_D$ is the repetition ratio for the entire dataset $D$, and each $S \in D$ is a sequence (LLM response) in $D$. The \rap score of the entire dataset is computed as follows:
\begin{align}
\label{eq:rap}
    RAP = P \times F(\rr)
\end{align}
}

\noindent
where $P$ can be any evaluation metric such as F1, BERT-F1, or COMET, and \new{$F(\rr)$ is the penalty function that calculates the penalty volume based on \rr. Unless otherwise specified, the results in this paper are generated using cubic penalty function, i.e., $F(\rr)=(1 - \rr)^3$.} 
If no repetition is found in the generated text, $RAP = P$; otherwise, $RAP < P$, meaning the performance is penalized for having repetition. \rap enables tuning \rp to achieve the balance between model performance and repetition, allowing for high performance while minimizing repetitive content.

% To compute the \rr of the whole dataset, we can take the average of \rr 
% compute \nrr of each LLM response. We then compute the mean value across the dataset to determine the repetition factor (\rf), defined as follows:
% \begin{align}
%     RF = \log_{10}(10 + \alpha \times MTP)
% \end{align}
% \noindent
% where  $MTP$  denotes the mean total repetitions, which is the average number of repetitions, including both \nwc and text repetitions, in each response. 
% If no repetition is found in the generated text, then  $RF = 1$; otherwise,  $RF > 1$ . The term  $\alpha$  is introduced to control the level of tolerance for repetition: the smaller the value of $\alpha$, the less penalization of repetition. In our study, we set  $\alpha = 1$.
% Here, $RF$ represents the repetition factor and $MTP$ denotes the mean total repetitions, i.e., the average repetition including \whitespace and text repetition in each response. If no repetition is found in the generated text, then $RF = 1$; otherwise, \( RF > 1 \). The term $\alpha$ is introduced to control the level of \textit{repetition tolerate}, the smaller value of $\alpha$, the less penalization of repetition. In our study, we set $\alpha = 1$ for simplicity.

% where $P$ can be any performance metric such as F1, BLEU, or ROUGE. If no repetition is found in the generated text, $RAP = P$; otherwise, $RAP < P$.
% \rap enables tuning \rp to obtain the trade-off between model performance and repetition that maintain high performance while reducing repetition content. 
% Instead of solely optimizing for the best performance, we can search for the best \( RAP \) to balance reducing excessive repetition while maintaining high performance.



